\section{Técnica de conteo}
\subsection{Notación factorial}
Se llama factorial a la operación de multiplicar todos los numeros enteros anteriores al número n, por definicón $0!=1$.
\begin{equation*}
    n!=n(n-1)(n-2)(n-3)\dots(1)(0!)
\end{equation*}
\subsection{Permutaciones}
La manera de ordenar un conjunto de objetos se le conoce como permutación. El número de permutaciones de $n$ objetos diferentes tomando $r$ al mismo tiempo se calcula como:
\begin{equation*}
    P\left(\begin{matrix}
            n \\
            r\end{matrix}\right) = \frac{n!}{(n-r)!} \qquad \text{donde} \qquad n\geq r
\end{equation*}
\begin{ejemplo}
    Encontrar el número total de permutaciones del conjunto de letras $\{a,b,c,d\}$ tomadas:
    \begin{itemize}
        \item tres a la vez
              \begin{align*}
                  P\left(\begin{matrix}
                          4 \\
                          3
                      \end{matrix}\right) & = \frac{4!}{(4-3)!} \\
                                                          & = \frac{4!}{1!}     \\
                                                          & = 4!                \\
                                                          & = 24
              \end{align*}
        \item dos a la vez
              \begin{align*}
                  P\left(\begin{matrix}
                          4 \\
                          2
                      \end{matrix}\right) & = \frac{4!}{(4-2)!}   \\
                                                          & = \frac{4!}{2!}       \\
                                                          & = \frac{4(3)(2!)}{2!} \\
                                                          & = 4(3)                \\
                                                          & =12
              \end{align*}
    \end{itemize}
\end{ejemplo}
\begin{teorema}
    El número de permutaciones diferentes de n objetos de los cualess $n_1$ son iguales, $n_2$ son iguales, $\dots$, $n_k$ son iguales es:
    \begin{equation*}
        P\left(\begin{matrix}
                n \\
                n_1,n_2,\dots,n_k
            \end{matrix}\right) = \frac{n!}{n_1!n_2!\dots n_k!}
    \end{equation*}
\end{teorema}
\begin{ejemplo}
    ¿De cuantas maneras distintas pueden situarse diez sillas en fila, si, cuatro son negras tres color verde y tres rojas?
    \begin{align*}
        P\left(\begin{matrix}
                10 \\
                4,3,3
            \end{matrix}\right) & = \frac{10!}{4!3!3!}             \\
                                                & = \frac{10(9)(8)(7)(6)(5)}{3!3!} \\
                                                & = \frac{10(9)(8)(7)(5)}{3!}      \\
                                                & = 10(3)(4)(7)(5)                 \\
                                                & =4200
    \end{align*}
\end{ejemplo}
\begin{teorema}
    El número de permutaciones de n objetos distitnos en un circulo es $(n-1)!$
\end{teorema}
\subsection{Combinaciones}
Una seleccion de n objetos en tomando r, sin tener en cuenta el orden, se le conoce como combinación y se calcula de la siguiente manera:
\begin{equation*}
    C\left(\begin{matrix}
            n \\
            r
        \end{matrix}\right) = \frac{P\left(\begin{matrix}
            n \\
            r
        \end{matrix}\right)}{r!} = \frac{n!}{r!(n-r)!}
\end{equation*}